\documentclass[openany,fontset=fandol]{ctexbook}
\usepackage{amsmath,amssymb}
\usepackage[unicode,hidelinks]{hyperref}
\newcommand{\note}[1]{\footnote{译注：#1}}
\begin{document}
\title{代数几何基础第五卷预备笔记\\引理13.3：超平面情形的证明文本}
\author{据用户上传英译编辑稿译出}
\date{2026年11月17日}
\maketitle
\chapter*{范围与校订状态}
本批连续译出底本第52页从“It remains to prove the corollary announced of 8.8 in the case $m=1$”起，至“this proves our regret, thus 13.2”止的一整段。为便于阅读，按论证步骤分段。下一处是同页“N.B. The reasoning does not go through”。

至此已译入引理13.3证明及定理13.2证明的原文收尾，但不能据此说数学证明已经独立验证。所引8.8的缺文、8.7(c)的残缺、一般平展条件及对偶论证的适用范围仍待核。下文校订两处连接词与字母错置，保留原写法及理由；没有将本段倒填为8.8的完整原文。本批未并入原合订本。
\chapter*{§13 初等态射与M. Artin定理（续）}
还须证明上面宣布的8.8的推论，即$m=1$的情形。如果$\operatorname{codim}T\geq2$，另一方面，参数化满足$x\in H_\xi$的超平面的子空间$Q$是对偶射影空间$P^\vee$中的一个超平面，余维为1；因此$Q$的一般点$\eta$不可能属于$T$，于是便得到了所需的结论。\note{底本在结句保留法语“on gagne [Fr]”，随后附有英译疑问“(we are done?)”，此处译出其完成论证之意，并保存该疑问。这里用到的是$T$作为闭的坏参数集所具有的含义：若$Q$的一般点属于$T$，则整个$Q$包含于$T$，与余维大小冲突；本批没有恢复8.8缺失的完整假设。}

在$\operatorname{codim}T=1$的情形，由于$T$不可约，只有当$Q=T$时，才可能有$\eta\in T$。\note{底本原句为“we cannot have $\eta\in T$ under $Q=T$”，按字面与此处论证相反。由$\eta\in T$得$Q\subseteq T$，而二者为同维不可约闭子簇，所以必须相等；后文也明确在$T=Q$的情形下继续论证。正文据此订为“只有当……时才可能”，原连接词保存在本注。}
用几何语言来说（可以假设$k$代数闭），这意味着：对每个$z\in X$，$X$在$z$处的切空间，确切地说，是该切空间经$f'_z$所得的像，经过$x$。\note{底本的代数闭括注没有正常闭合；这里整理括号，不补造缺失论证。原“or also plutot [Fr]”按法语plutôt译为“确切地说”。切空间在此按原射影对偶论证理解，不将“经过$x$”当成抽象向量空间的无坐标陈述。原段没有逐项重述对偶簇$T$及切空间像的定义，适用条件仍须与前文核对。}
\newpage
我们来证明，在$X^{\mathrm{sing}}\to T$一般平展、也就是处于所引用的8.8条件之下时，这种情形不能发生，除非$f(X)=\{x\}$，此时$X$为零维。\note{底本两次使用unless：“cannot happen unless $X^{\mathrm{sing}}\to T$ is generically étale, i.e. unless we are under the conditions of 8.8”。但前一批原段把一般平展列为获得结论的充分条件，紧随其后的论证也正是在该条件下引用8.7(c)得出矛盾；正文依此订为“在……条件下不能发生”。这一校订只处理连接词与上下文的逻辑方向，不补齐8.8的缺失条件，也不宣称该对偶结论已获独立证明。此处原印$X^{\mathrm{sing}}$，上一页印$h^{\mathrm{sing}}$，记号不一致照录待核；不能未经核对就将它解释为当前光滑$X$的奇异点集。底本$f(x)=\{x\}$据同段对所有$z$及$X$维数的讨论订为$f(X)=\{x\}$。在原先$f$未分歧的有限型背景中，像为一点对应零维情形；并非声称任意常值态射的定义域都零维。}

事实上，由8.7(c)（它本质上表达了$X$与其“对偶”$T$之间关系的对称性），对于几乎所有$z\in X(k)$，$f(z)$与$T$在某一点处的切空间正交。由于$T=Q$，这就意味着$f(z)$与$Q$正交，故$f(z)=x$，从而$f(X)=\{x\}$。这证明了我们的补记，因而也证明了13.2。\note{底本后半句语序破碎：“at a certain point this (since $T=Q$) orthogonal to $Q$, where $f(z)=x$”；正文据其显式等式整理为连续推论。在射影对偶意义下，$Q$是所有经过$x$的超平面的参数超平面，其正交点就是$x$。末句保留的是原作者关于论证完成的表述，不是译者对全部依赖链的独立审定。“几乎所有”在此是代数几何中一般点意义的说法，不是概率或测度断言；从一般$z$的结论推广到整个像所需的稠密性及闭性使用在原稿中未展开。尤其8.7(c)和8.8的原文残缺仍未由本段恢复。}
\end{document}
